\chapter{Dal perceptron alla rete neurale}
\label{cap:reti-neurali}
\fonti{\rif{2026-09-28}{28/09}}{2-foundament-deep-learning, pp.\ 1--18}

\section{Il quadro generale}
\begin{itemize}
  \item \textbf{Input layer}: riceve i dati. Con dati tabellari, ogni riga è un punto e ogni colonna una feature
        ($x_1, x_2, \ldots$): un neurone di input per feature.
  \item \textbf{Hidden layer}: la parte che ``ragiona'' e modella il problema. Con 3 o più hidden layer si parla di
        \textbf{deep learning}.
  \item \textbf{Output layer}: prende la decisione finale (es.\ cane, gatto o delfino).
\end{itemize}

\begin{attenzione}[le reti mangiano solo numeri]
Testo, grafi e immagini vanno convertiti in rappresentazioni numeriche. Un'immagine è una griglia di pixel e
ogni pixel è descritto da 3 numeri (RGB). Questa conversione è di solito standard; le parti difficili sono:
\begin{itemize}
  \item il \textbf{preprocessing} che rende i dati puliti e utili;
  \item l'\textbf{architettura}: lo stesso problema può essere visto come tabella (FFNN), serie temporale (RNN) o grafo (GNN),
        e va scoperto sperimentalmente quale funziona meglio, quanti layer e quanti neuroni usare
        (\textbf{model engineering}).
\end{itemize}
\end{attenzione}

\section{Il perceptron}
\begin{definizione}[perceptron]
L'unità più semplice di una rete neurale. Prende un input con più feature e:
\begin{enumerate}
  \item pesa ogni feature con un \textbf{peso} $w_i$;
  \item calcola la \textbf{somma pesata};
  \item applica una \textbf{funzione di attivazione} $f(\cdot)$ e produce l'output:
\end{enumerate}
\[
y = f\Big(\sum_{i=0}^{n} w_i x_i\Big) = f(\mathbf{w}^T \mathbf{x}), \qquad x_0 = 1 \text{ per definizione}.
\]
\end{definizione}

\begin{itemize}
  \item I \textbf{pesi} sono i \textbf{parametri} della rete (i ``trilioni di parametri'' degli LLM): indicano
        l'\textbf{importanza} di ogni feature per il compito. All'inizio sono \textbf{casuali} e vengono
        \textbf{appresi dai dati}.
  \item Le $x_i$ non si toccano: sono i valori delle feature, dati dal mondo esterno. Si impara solo $\mathbf{w}$.
  \item Il \textbf{bias} $w_0$ è un peso come gli altri ma \textbf{non dipende da una feature} (moltiplica $x_0=1$):
        serve a \textbf{traslare} la funzione.
  \item Nei layer successivi l'input non sono le feature ma gli \textbf{output del layer precedente}; la struttura del neurone
        è identica.
\end{itemize}

\begin{esempio}[perché serve il bias]
Con attivazione lineare ($f(x)=x$) e un solo input, $y = w_1 x_1$: posso cambiare solo la \textbf{pendenza}, e tutte le
rette passano per l'origine. Se il confine tra croci e punti non passa per l'origine, non posso modellarlo. Aggiungendo
$w_0$ ottengo $y = w_1x_1 + w_0$ e posso spostare la retta.
\end{esempio}

\begin{integrazione}[legame con modelli classici]
Senza $f(\cdot)$ il perceptron è una \textbf{regressione lineare}; con $f = \sigma$ (sigmoid) è una
\textbf{regressione logistica}.
\end{integrazione}

\subsection{Caso 2D con attivazione lineare}
$y = w_1x_1 + w_2x_2 + w_0$ descrive una famiglia di funzioni (un piano/una retta di separazione).
\begin{itemize}
  \item Se tutti i pesi sono 0, l'output è sempre 0: la rete non predice nulla.
  \item Tarando $w_1$ separo lo spazio lungo $x_1$; con $w_2$ inclino il confine; con $w_0$ lo sposto.
  \item L'output grezzo non è limitato: con $x_2=10$ può valere molto più di 10. Per avere una
        \textbf{probabilità} serve una funzione limitata in $[0,1]$.
\end{itemize}

\section{Funzioni di attivazione}
Servono a:
\begin{enumerate}
  \item \textbf{imporre proprietà} all'output (es.\ sigmoid $\Rightarrow$ output in $[0,1]$);
  \item introdurre \textbf{non linearità} nel modello;
  \item decidere \textbf{quando il neurone è attivo} (ReLU: se l'input è $\leq 0$ il neurone è spento);
  \item altri motivi di ottimizzazione: convergenza più rapida, \en{sparsity}.
\end{enumerate}

\begin{center}
\small
\renewcommand{\arraystretch}{1.35}
\begin{tabularx}{\textwidth}{l l l X}
\toprule
\textbf{Funzione} & \textbf{Formula} & \textbf{Codominio} & \textbf{Uso tipico e note}\\
\midrule
Sigmoid & $\sigma(y)=\dfrac{1}{1+e^{-y}}$ & $(0,1)$ & output di classificazione \textbf{binaria}; limitata $\Rightarrow$ \en{vanishing gradient} se usata negli hidden layer\\
Tanh & $\tanh(y)$ & $(-1,1)$ & limitata; usata nelle \textbf{RNN}\\
ReLU & $\max(0,y)$ & $[0,\infty)$ & la più usata negli \textbf{hidden layer}; non limitata superiormente; neuroni inattivi $\Rightarrow$ sparsità\\
Leaky ReLU & $y$ se $y>0$, $\alpha y$ altrimenti & $\R$ & piccola pendenza per $y<0$: anche i neuroni ``negativi'' portano informazione\\
Softmax & $\dfrac{e^{z_i}}{\sum_j e^{z_j}}$ & $(0,1)$, somma 1 & output \textbf{multi-classe}\\
Lineare & $f(y)=y$ & $\R$ & output di \textbf{regressione}\\
\bottomrule
\end{tabularx}
\end{center}

\subsection{Sigmoid}
Mappa qualsiasi valore reale in $[0,1]$: anche se la somma pesata vale $-10.000$, l'output è limitato.
$P(\text{pos}\mid x) = \sigma(y)$ e $P(\text{neg}\mid x) = 1 - P(\text{pos}\mid x)$, quindi la somma è 1 per costruzione.
Soglia tipica: $0{,}5$. Con tutti i pesi a 0 l'output è $\sigma(0)=0{,}5$: il modello non sa decidere.

\textbf{Limiti della sigmoid}
\begin{enumerate}
  \item \textbf{Solo binaria}: per cane/gatto/delfino non basta.
  \item \textbf{Vanishing gradient}: in addestramento si manda indietro (\en{backpropagation}) un segnale che dipende dalle
        derivate delle attivazioni. Moltiplicando valori $<1$ per molti layer, il segnale \textbf{svanisce} e i primi layer non
        imparano (dettagli nel cap.~\ref{cap:addestramento}).
\end{enumerate}

\subsection{ReLU}
Se l'input è $\leq 0$ l'output è 0 (neurone \textbf{inattivo} per quella combinazione di feature); altrimenti l'output è
proporzionale all'input. Non essendo limitata, il segnale backpropagato resta abbastanza forte anche dopo molti layer.
Benefici: mitiga il vanishing gradient $\Rightarrow$ \textbf{convergenza più rapida}; \textbf{sparsità}: non tutti i neuroni
servono per tutti gli input.

\begin{esempio}[sparsità e LLM]
Un LLM da trilioni di parametri non li usa tutti a ogni richiesta: alcuni neuroni sono specializzati (es.\ ``inglese'' più
che ``francese''). È collegato all'idea di \en{mixture of experts} e di \en{routing} dell'informazione nella rete.
\end{esempio}

\begin{attenzione}[ReLU non risolve tutto]
Moltiplicando molti valori $>1$ il gradiente può \textbf{esplodere} (\en{exploding gradient}): il problema opposto.
Succede anche con reti non troppo profonde. Reti diverse usano attivazioni diverse o combinazioni di attivazioni.
\end{attenzione}

\section{Attivazioni in uscita e logit}
\begin{definizione}[logit]
L'output della rete \textbf{prima} dell'attivazione finale: valori non normalizzati in $(-\infty, +\infty)$.
Vanno trasformati in probabilità valide (in $[0,1]$ e con somma 1).
\end{definizione}

\subsection{Classificazione binaria}
La rete produce \textbf{un logit} $z$ (per convenzione della classe positiva). Si applica la sigmoid:
$P(\text{pos}\mid x) = \sigma(z) = \frac{1}{1+e^{-z}}$ e $P(\text{neg}\mid x) = 1 - P(\text{pos}\mid x)$.
Con logit $-10.000$ la probabilità di positivo è $\approx 0$; con $+10.000$ è $\approx 1$.

\subsection{Classificazione multi-classe}
Servono almeno $C-1$ output, ma in pratica si usa \textbf{un neurone per classe} ($C$ logit). Si applica la
\textbf{softmax}:
\[
P(c_i \mid x) = \frac{e^{z_i}}{\sum_{j} e^{z_j}}, \qquad P(c_i\mid x) \in [0,1], \qquad \sum_i P(c_i \mid x) = 1.
\]
\begin{itemize}
  \item È una generalizzazione della sigmoid e usa \textbf{tutti gli output insieme}.
  \item Converte i logit in probabilità e \textbf{amplifica} le piccole differenze tra logit.
  \item Esempio: immagine $\to$ CNN (rappresentazione) $\to$ FFNN (classificazione) $\to$ 4 logit $\to$ softmax $\to$
        \{cane, gatto, uccello, tartaruga\}; si sceglie la probabilità più alta.
\end{itemize}

\subsection{Regressione}
Si predicono numeri reali: di solito non serve imporre proprietà e l'attivazione di uscita è l'\textbf{identità}
$f(x)=x$. È praticamente l'unico caso in cui ha senso un'attivazione lineare.

\begin{attenzione}[teoria vs software]
In teoria in classificazione l'ultima attivazione c'è sempre. In \textbf{PyTorch} però, con \code{CrossEntropyLoss},
\textbf{non} si mette la softmax nel modello perché è già dentro la loss (cap.~\ref{cap:ffnn}). È una scelta della libreria:
leggi sempre la documentazione.
\end{attenzione}
\section{Il margine: quale confine è migliore?}
Tra due classificatori lineari che separano correttamente i punti di training, è preferibile quello con
\textbf{margine maggiore} dai punti: un nuovo punto vicino al confine viene classificato correttamente con più probabilità.
Un confine troppo ``attaccato'' a una classe sbaglierà facilmente nei dintorni. È l'idea delle \textbf{SVM}
(cap.~\ref{cap:anomaly}).

\section{Verso reti più grandi}
\subsection{Impilare perceptron: il layer fully connected}
Due perceptron con lo stesso input $(x_1,x_2)$ e pesi diversi ($\mathbf{w}$ e $\mathbf{q}$) danno due output indipendenti:
\[
\mathbf{y} = \begin{pmatrix} y_1\\ y_2\end{pmatrix}
= f\left(\begin{pmatrix} w_0 & w_1 & w_2\\ q_0 & q_1 & q_2\end{pmatrix}
\begin{pmatrix} 1\\ x_1\\ x_2\end{pmatrix}\right) = f(W\mathbf{x})
\]
Questo è un \textbf{layer lineare (fully connected)} con \textbf{6 parametri}. Per confronto, i modelli ``piccoli''
di oggi ne hanno 7 miliardi (es.\ Mistral 7B).

\begin{definizione}[numero di parametri]
È il numero totale di \textbf{pesi} (bias inclusi). Un layer lineare da $n_{in}$ a $n_{out}$ neuroni ha
$n_{in} \cdot n_{out} + n_{out}$ parametri.
\end{definizione}

\subsection{Aggiungere layer}
Un secondo layer combina $y_1, y_2$ con pesi $\mathbf{s}$ e produce $z$:
\[
z = f\big(\mathbf{s}^T f(W^T\mathbf{x})\big)
\]
\textbf{Una rete neurale è moltiplicazione di matrici} alternata a funzioni di attivazione: nessuna magia, ma con
trilioni di pesi.

\textbf{Perché tanti layer?} Ogni layer crea una rappresentazione \textbf{più astratta} dei dati: i primi layer
lavorano sul dato grezzo, i successivi catturano i concetti principali (sarà evidentissimo con le CNN).

\textbf{L'attivazione è la stessa in tutti i layer?} No. Negli hidden layer serve evitare il vanishing gradient
$\Rightarrow$ ReLU (non limitata); nell'ultimo layer serve una probabilità $\Rightarrow$ sigmoid o softmax.

\section{Perché la non linearità è indispensabile}
Se $f(x)=x$, allora $z = \mathbf{s}^T W^T \mathbf{x} = W'\mathbf{x}$:
\begin{enumerate}
  \item potevo usare direttamente $W' = W\mathbf{s}$ e ottenere lo stesso output;
  \item il secondo layer non serviva;
  \item il modello resta \textbf{lineare}.
\end{enumerate}
Mille layer lineari equivalgono a \textbf{un solo layer lineare} (\en{collapse}). Il mondo non è lineare (pensa alla
distribuzione dei voti di un esame), quindi servono le attivazioni non lineari tra un layer e l'altro.

Con ReLU nella rete del docente (hidden ReLU + output sigmoid) il confine di decisione non è più una retta: si formano
\textbf{più regioni}. Il numero di regioni dipende da quanti neuroni ci sono e dalla presenza della non linearità.

\section{Teorema di approssimazione universale}
\begin{definizione}[Cybenko 1989; Hornik et al.\ 1989]
Per ogni funzione continua $g$ definita su un sottoinsieme compatto di $\R^n$ e per ogni $\epsilon > 0$, esiste una
rete feed-forward con \textbf{un solo hidden layer} e un numero finito di neuroni che approssima $g$ con errore inferiore a $\epsilon$.
\end{definizione}
\begin{itemize}
  \item Ogni neurone ReLU introduce un \textbf{punto di rottura} (\en{breakpoint}): la funzione approssimata è lineare a tratti.
        Con un input e $D$ neuroni si ottengono fino a $D+1$ \textbf{regioni lineari}. Nell'esempio delle slide, la curva viene
        approssimata sempre meglio con 4, 9 e infine \textbf{19 neuroni} ($\Rightarrow$ 20 regioni lineari): una rete
        \en{shallow} basta.
  \item Il numero di neuroni di una rete shallow misura la sua \textbf{capacità}.
  \item Il problema: il teorema \textbf{non dice quanti neuroni servono} (4? 19? 20.000?) né quali pesi usare.
\end{itemize}

\begin{esame}[meglio profonde che larghe]
Si usano reti \textbf{più profonde e più strette} invece di un solo layer larghissimo: i layer, separati da non linearità,
\textbf{combinano la loro complessità} e creano rappresentazioni sempre più astratte $\Rightarrow$ migliore
\textbf{generalizzazione}. ``Il deep learning è rappresentazione dei dati.''\\
Attenzione però: una rete troppo complessa può andare in \textbf{overfitting} (cap.~\ref{cap:addestramento}).
\end{esame}

\subsection{Complessità di una rete}
Il numero di regioni lineari cresce rapidamente con il numero di neuroni e di \textbf{feature di input}.
Con 500 neuroni e 100 feature si superano le $10^{107}$ regioni, con \textbf{51.001 parametri}: tantissime regioni, eppure
la rete è minuscola per gli standard attuali (Mistral 7B ha 7 miliardi di parametri).


\begin{sintesi}
\begin{itemize}
  \item Il perceptron calcola $y=f(\mathbf{w}^T\mathbf{x})$: pesi appresi dai dati, bias $w_0$ per traslare il confine,
        attivazione $f$ per imporre proprietà e non linearità.
  \item Attivazioni: ReLU negli hidden layer (mitiga il vanishing gradient, sparsità); sigmoid (binaria), softmax
        (multi-classe) o identità (regressione) in uscita. I \textbf{logit} sono gli output prima dell'attivazione finale.
  \item Un layer fully connected è $f(W\mathbf{x})$, con $n_{in}n_{out}+n_{out}$ parametri; una rete è moltiplicazione di
        matrici alternata a non linearità. Senza non linearità, mille layer collassano in uno solo.
  \item Teorema di approssimazione universale: basta un hidden layer, ma non dice quanti neuroni né quali pesi.
        In pratica si preferiscono reti \textbf{profonde e strette}, che generalizzano meglio.
\end{itemize}
\end{sintesi}
